% s00-front.tex —— 译者说明
\chapter*{译者说明}
\addcontentsline{toc}{chapter}{译者说明}
\markboth{译者说明}{译者说明}

\section*{底本}

Edward Frenkel、Robert Langlands、Ngô Bao Châu，
\emph{Formule des traces et fonctorialité : le début d'un programme}，
arXiv:1003.4578v2 [math.RT]，2010 年 9 月 29 日，54 页。
原文为\textbf{法语}。印刷页码与 PDF 页码一致（偏移 0）：正文 p.1--52，文献 p.53--54。

\section*{译本体例}

\begin{enumerate}[leftmargin=2em,itemsep=0.2em]
  \item \textbf{节号、公式号、断言号一律照原文}。原文的 \emph{Lemme 2.6}、\emph{Proposition 3.29} 等，
        译本作「引理 2.6」「命题 3.29」，编号不变，可直接与原文对照。
        原文未编号的节（导引）译本亦不编号。
  \item \textbf{页边 } \texttt{p.$N$} \textbf{ 标记原文页码}，N 为底本印刷页。
        引用、校勘一律用原书页。
  \item \textbf{译注}以脚注排出，楷体，末尾用 \texttt{[来源：…]} 给出依据：
        A 级为原书、作者论文与经典教材；B 级为 nLab、Stacks Project 等；
        I 级为本库既有译注与术语表。凡无外部依据、仅由上下文推断者，
        标 \texttt{[译者按（无外部依据）]}。
  \item \textbf{术语}按 \texttt{工程文件/work/glossary.tsv} 冻结；首次出现处酌情附原文。
        有通行中文译名者用中文；无通行译名或多译法者保留原文。
  \item \textbf{原文的笔误、记号的混淆}一律保留正文原样，另加译注说明，不擅自改动。
  \item 方括号中的 \texttt{[N]}、\texttt{[L3]} 等为原文文献编号，译本保留原编号与文献表顺序。
\end{enumerate}

\section*{关于译文}

本译本是数译工程（AI 辅助）的产物：逐页对照原文与页面图翻译，
公式逐符号核对，术语与文献经本库知识库交叉检查。
数学内容以原文为准；读者若发现译文与原文不符，请以原文为据，并欢迎指出。
